%-------------------------
% CV Conforme ATS — Omar Bouhlal (Français)
% Utilise uniquement du LaTeX standard — compatibilité ATS maximale
%-------------------------
\documentclass[a4paper,11pt]{article}

% --- Packages ---
\usepackage[utf8]{inputenc}
\usepackage[T1]{fontenc}
\usepackage{lmodern}
\usepackage[margin=0.5in]{geometry}
\usepackage[inline]{enumitem}
\usepackage{hyperref}
\usepackage{titlesec}
\usepackage{xcolor}
\usepackage[french]{babel}

% --- Colors ---
\definecolor{accent}{HTML}{2563EB}
\definecolor{darktext}{HTML}{1f2937}
\definecolor{lighttext}{HTML}{6b7280}

% --- Hyperlinks ---
\hypersetup{
    colorlinks=true,
    linkcolor=accent,
    urlcolor=accent,
    pdfauthor={Omar Bouhlal},
    pdftitle={Omar Bouhlal - CV},
}

% --- Section Formatting ---
\titleformat{\section}
  {\large\bfseries\color{darktext}\uppercase}
  {}{0em}{}
  [\titlerule]
\titlespacing*{\section}{0pt}{10pt}{4pt}

% --- No page numbers ---
\pagestyle{empty}

% --- Custom commands ---
\newcommand{\entryhead}[3]{%
  \textbf{#1} \hfill \textit{#2} \\
  \textit{#3}
}

\begin{document}

% ============================================================
% EN-TÊTE
% ============================================================
\begin{center}
  {\LARGE\bfseries Omar Bouhlal} \\[2pt]
  {\color{lighttext} Meknès, Maroc} \\[2pt]
  omarbouhlal05@gmail.com
  \quad $\cdot$ \quad +212 766 678 608 \\[2pt]
  linkedin.com/in/omar-bouhlal-101ba6353
  \quad $\cdot$ \quad
  github.com/OmarBouhlal
\end{center}

\vspace{-10pt}



% ============================================================
% FORMATION
% ============================================================
\section{Formation}

\entryhead{ENSAM Meknès — École Nationale Supérieure d'Arts et Métiers}{2022 -- Présent}{Diplôme d'Ingénieur en Génie Logiciel et Systèmes Intelligents (4ème année)}

\vspace{2pt}

\entryhead{Groupe Scolaire Atlas — Rabat}{2021 -- 2022}{2ème Année Baccalauréat Sciences Physiques — Mention Très Bien}

\vspace{2pt}

\entryhead{Groupe Scolaire Atlas — Rabat}{2020 -- 2021}{1ère Année Baccalauréat Sciences Expérimentales — Mention Très Bien}

% ============================================================
% EXPÉRIENCE PROFESSIONNELLE
% ============================================================
\section{Expérience Professionnelle}

\entryhead{ECS Informatique}{2024 · 1 mois}{Stagiaire en Ingénierie Logicielle}
\begin{itemize}[leftmargin=*, nosep, topsep=2pt]
  \item Développement d'une application web d'entreprise pour la gestion et le suivi multicritère des frais professionnels des employés
  \item Conception et implémentation de formulaires dynamiques, contrôle d'accès par rôles et tableaux de bord de reporting
  \item Collaboration en équipe sur la conception de l'architecture de la base de données et des API REST
  \item Technologies : ASP.NET, C\#, .NET Framework, SQL Server
\end{itemize}

% ============================================================
% PROJETS
% ============================================================
\section{Projets}

\textbf{ENSAM360 — Visite Virtuelle Immersive 360°} \\
\textit{React, Node.js, Neo4j, Groq AI, Docker}
\begin{itemize}[leftmargin=*, nosep, topsep=2pt]
  \item Conception d'une plateforme de visite virtuelle permettant la navigation numérique du campus de l'ENSAM Meknès
  \item Intégration d'un chatbot alimenté par l'IA pour assister les visiteurs durant leur expérience de visite
  \item Architecture micro-services orchestrée via Docker Compose
\end{itemize}

\vspace{4pt}

\textbf{XPFit — Application Mobile de Fitness} \\
\textit{Flutter, Dart, SQLite, Google AI, Spoonacular API}
\begin{itemize}[leftmargin=*, nosep, topsep=2pt]
  \item Développement d'une application mobile multiplateforme avec suivi d'exercices gamifié et système d'évolution physique
  \item Implémentation de plans de repas hebdomadaires avec guide nutritionnel via l'API Spoonacular
  \item Intégration d'un chatbot IA utilisant Google Generative AI pour des conseils fitness personnalisés
\end{itemize}

\vspace{4pt}

\textbf{Jeu d'Échecs — Pixel Art 2D} \\
\textit{Godot Engine, GDScript, LibreSprite}
\begin{itemize}[leftmargin=*, nosep, topsep=2pt]
  \item Création d'un jeu d'échecs 2D entièrement jouable à partir de zéro avec le moteur Godot
  \item Réalisation de tous les sprites des pièces en pixel art avec LibreSprite
\end{itemize}

% ============================================================
% COMPÉTENCES TECHNIQUES
% ============================================================
\section{Compétences Techniques}

\textbf{Langages :} Java, C\#, Dart, GDScript, JavaScript, Python, C++ \\
\textbf{Frameworks :} Spring Boot, ASP.NET / .NET, Flutter, React, Node.js \\
\textbf{Bases de données :} SQL Server, SQLite, Neo4j \\
\textbf{Outils et Plateformes :} Git, Docker, Godot Engine, LibreSprite \\
\textbf{Concepts :} Microservices, API REST, Architecture Logicielle, Design Patterns

% ============================================================
% CERTIFICATIONS
% ============================================================
\section{Certifications}

\textbf{Oracle Cloud Infrastructure (OCI)} \hfill Oracle

% ============================================================
% LANGUES
% ============================================================
\section{Langues}

\textbf{Arabe} — Langue maternelle \quad $\cdot$ \quad
\textbf{Français} — Courant \quad $\cdot$ \quad
\textbf{Anglais} — Courant

\end{document}
